\documentclass[11pt]{article}
\usepackage[margin=1in]{geometry}
\usepackage[T1]{fontenc}
\usepackage{tgpagella,microtype,booktabs,amsmath,graphicx,pdflscape}
\usepackage[hidelinks]{hyperref}
\title{Forecasting with a Fixed Evidence Set:\\A Three-Model Comparison against Prediction-Market Probabilities}
\author{Vorhersage experimental series}
\date{September 18, 2026}
\begin{document}
\maketitle
\begin{abstract}
We compare three frontier models on six prediction questions, holding the question
definitions, evidence packets, elicitation prompts, and market benchmarks fixed.
Thirty-six new calls cross three models with question-only and outside-research
conditions. All eighteen research forecasts pass validation. Their mean absolute
distances from the original Kalshi midpoints are 5.42 percentage points for Claude
Fable 5.1, 6.25 for GPT-6 Astra, and 8.42 for Gemini 3.1 Pro. The original Gemini
run yielded 8.58 points. Research reduces average distance for every model on the
three contracts with accepted responses in all six arms. These observations
illustrate a controlled comparison of forecasting methods using timely market
benchmarks; the small sample and single draw per cell do not establish a general
ranking of forecasting ability.
\end{abstract}

\section{Design}
The original study selected six contracts spanning weather, employment, monetary
policy, space technology, entertainment, and baseball. Its outside evidence came
from primary sources rather than Kalshi or other prediction markets. This
follow-up retains those exact questions and evidence packets. Three models each
receive the six questions twice: once without current research and once with the
corresponding frozen packet. The calls are independent; a research-condition call
does not see a question-only forecast. All 36 returned system and user prompt
receipts exactly match their respective prompts from the original study.

Prediction-market probabilities provide informative intermediate evaluation
targets and immediate feedback while event outcomes remain unavailable. Evidence
for their forecasting value and calibration, subject to liquidity and horizon,
is discussed in the companion report. We measure agreement with the same opening
midpoints used there, without downloading new prices or outcomes.

The information cutoff remains September 18, 2026, 11:19 UTC. New inference
occurred later that evening Eastern time (September 19 UTC). The coordinator had
already seen the original prices and forecasts; new model contexts contained
neither. These are therefore registered, frozen-data follow-up runs rather than
a new fully blind prospective study. The same-day weather outcome may have become
known by the later execution time. The original advantage of evaluating genuinely
future events before their outcomes exist should not be attributed unconditionally
to that later weather call.

\section{Evaluation}
Let $p_{imc}$ denote the forecast for question $i$, model $m$, and information
condition $c$, and let $q_i$ denote the original market midpoint. On a common
eligible cohort $S$, the principal statistic is
\[
 \operatorname{MAE}_{mc}=\frac{100}{|S|}\sum_{i\in S}|p_{imc}-q_i|.
\]
We also report root mean squared distance and preserve distances outside each
original bid--ask spread. Research improvement for an individual question is
$100(|p_{im0}-q_i|-|p_{im1}-q_i|)$ percentage points. Its relative improvement
divides that change by the original absolute distance. A negative value denotes
deterioration. These are market-agreement statistics, distinct from outcome
scores. Under the stronger assumption that $q_i$ is the conditional event
probability for the reference information set, squared distance also equals
conditional expected excess Brier loss; marginal calibration alone does not
establish that assumption.

\section{Results}
Table~\ref{tab:research} uses all six questions because every research forecast
passed validation. Claude is closest on this sample, followed by Astra and
Gemini. The new Gemini mean closely resembles its earlier mean, although
question-level estimates change. With only one new draw per cell, this is a
replication observation rather than a precise estimate of stochastic variability.
Excluding same-day weather leaves research MAE of 5.60 points for Claude,
6.20 for Astra, and 9.80 for Gemini on the same five remaining contracts.

\begin{table}[ht]
\centering
\caption{Research forecasts on the same six contracts. Distances are percentage
points. Cost includes both information conditions and rejected responses.}
\label{tab:research}
\begin{tabular}{lrrr}
\toprule Model & MAE & RMSE & Cost (USD) \\
\midrule
Claude Fable 5.1 & 5.42 & 6.99 & 0.84760 \\
GPT-6 Astra & 6.25 & 8.68 & 0.39430 \\
Gemini 3.1 Pro, new & 8.42 & 11.70 & 0.17177 \\
Gemini 3.1 Pro, original & 8.58 & 11.40 & 0.16835 \\
\bottomrule
\end{tabular}
\end{table}

The unchanged conservative text rule excluded three question-only responses:
Claude on employment and weather, and Gemini on Netflix. Each mentions prediction
markets and infers information from contract selection or threshold placement,
despite instructions treating contract rules as definitions only. Claude's
employment response also includes an explicit disclaimer of using odds; a mere
mention is not proof of price exposure. These responses remain excluded without
repair. Acceptance is 12/12 for Astra, 10/12 for Claude, and 11/12 for Gemini.

The fully matched before/after cohort consequently contains Fed, Neutron, and
baseball. Table~\ref{tab:paired} compares identical questions in all arms. Dropping
same-day weather, as prespecified, does not change this cohort. Available-case
results are archived separately rather than substituted for matched comparisons.

\begin{table}[ht]
\centering
\caption{Fully matched comparison on three contracts.}
\label{tab:paired}
\begin{tabular}{lrrr}
\toprule Model & Question-only MAE & Research MAE & Reduction \\
\midrule
GPT-6 Astra & 18.00 & 7.67 & 57.4\% \\
Claude Fable 5.1 & 17.00 & 4.00 & 76.5\% \\
Gemini 3.1 Pro, new & 30.33 & 8.00 & 73.6\% \\
\bottomrule
\end{tabular}
\end{table}

Figure~\ref{fig:slopes} preserves the question-level comparison. Effects are
heterogeneous: Claude's Netflix forecast moves slightly away from the midpoint,
and Gemini's Fed forecast does not move. The unchanged Netflix packet also
retains incomplete release context from the original study; replacing it would
constitute a different data experiment. A larger repeated comparison could
separate model variation from evidence-quality effects.

\clearpage
\begin{figure}[p]
\centering
\includegraphics[width=\textwidth,height=.91\textheight,keepaspectratio]{figures/performance.pdf}
\caption{Fixed-data model comparison. Slopes connect independent question-only
and research forecasts. Dashed lines mark the original Kalshi midpoints.
Annotations report absolute and relative gap reductions. Excluded responses have
no point or connecting segment; their accepted counterpart is retained.}
\label{fig:slopes}
\end{figure}
\clearpage
\appendix
\section{Implementation and reproducibility}
Vorhersage separates deterministic record keeping from external model judgment.
The experiment registers six immutable arms, each combining the same
direct-judgment method with a specific model and evidence condition. Frozen
packets carry source provenance and content-derived identifiers. Every trial
has independent assessment and issuance state; deterministic issuance requires
no additional model call. Runs are marked retrospective and preserve actual
execution times rather than backdating forecasts to the information cutoff.

Expected Parrot transports one fresh, tool-free EDSL interview per cell. Astra
uses high reasoning effort, Claude adaptive thinking at high effort, and Gemini
the original 2,048-token thinking budget and temperature 0.5. All have an
8,192-token output ceiling. These are model configurations rather than matched
computation budgets. Provider adapters may omit serialized legacy defaults;
the archive verifies returned model identities and prompt receipts without
claiming independent attestation of server-side settings or metering.

Import validates scenario, model, agent, JSON schema, evidence references, and
the excluded-content rule. An initial importer defect compared plain scenario
dictionaries with EDSL's serialized dictionaries and rejected all rows. The
corrected validator requires the two exact serialization metadata fields and
unchanged substantive content. Initial rejection records are retained. Neither
forecasts nor registered inputs changed, and no additional inference occurred.

Reported new response costs total \$1.41367, including excluded responses; all
costs are known in the returned records. Provider job summaries report twelve
completed interviews and no interviews with exceptions for each model. Source
collection and coordination costs are outside this total. Raw responses, model
settings, prompt hashes, submissions, acceptance records, issued forecasts,
and seal commitments are archived under \texttt{run/}.

The comparison script is \texttt{compare.py}. A separate ggplot2 script,
\texttt{plot\_performance.R}, generates the figure. Running \texttt{verify}, \texttt{audit},
and \texttt{score} commands rechecks inputs and reconstructs the comparison
without new inference or web access. The new boundary tests cover unchanged
prompts, serialization metadata, unexpected context, changed model settings,
and retention of the strict output exclusion policy. No dependency is added to
the core package.
\end{document}
